\documentclass[10pt,a4paper]{amsart}
\usepackage[margin=19mm]{geometry}
\usepackage{amsmath,amssymb,mathtools}
\usepackage[scr=rsfs,cal=euler]{mathalfa}
\usepackage{xfrac}
\usepackage[unicode,hidelinks,bookmarks=false]{hyperref}
\newcommand{\ie}{i.e.\ }
\newcommand{\cf}{cf.\ }
\newcommand{\resp}{respectively\ }
\newcommand{\Subsection}{\subsection}
\providecommand{\nobreakdash}{\nobreak\hbox{-}}
\setlength{\emergencystretch}{2em}
\multlinegap0pt
\usepackage{amssymb}
\usepackage{cancel}
\usepackage{caption}
\usepackage{tikz}
\newtheorem{example}{Example}
\newtheorem{lemma}{Lemma}
\DeclareMathOperator{\cn}{cn}
\DeclareMathOperator{\dn}{dn}
\DeclareMathOperator{\sn}{sn}
\title{\bf Bright and dark breathers on an elliptic wave \\in the defocusing mKdV equation}
\author{Dmitry E. Pelinovsky, Rudi Weikard}
\date{Source: arXiv:2512.02959v2 (CC BY 4.0)}
\begin{document}
\maketitle

\noindent\emph{Source and licence.} \bf Bright and dark breathers on an elliptic wave \\in the defocusing mKdV equation. Version arXiv:2512.02959v2 (CC BY 4.0), distributed by arXiv under the Creative Commons Attribution 4.0 International licence, http://creativecommons.org/licenses/by/4.0/, as recorded in the arXiv licence metadata for this exact version and shown on the abstract page https://arxiv.org/abs/2512.02959v2. Full licence notice: \texttt{licenses/arxiv-2512.02959-CC-BY-4.0.txt}.\par\smallskip

\noindent\textit{Editorial scope.} Counted unit: the article's lemma lem-coefficients (source0.tex lines 718--730). The article's complete original proof of it is delivered with it (source0.tex lines 733--870, one \texttt{proof} environment of 138 lines). No mathematics is written, repaired or re-derived: the statement and the proof are the article's own lines, in the article's own notation, and no retained line is substituted or re-flowed.  Retained uncounted as the source-local context the proof needs: lines 148--151; lines 227--231; lines 275--279; lines 281--285; lines 287--292; lines 300--303; lines 499--502; lines 561--564; lines 576--579; lines 602--607; lines 673--680; lines 682--688; lines 691--713.  One declared adaptation: exactly 1 ancillary strings inside the retained spans are omitted (image-inclusion commands and, where the source repeats a label, the repeated label definitions) and no other line differs from the source; the figure environments, with their placement, captions and labels, are kept, and the package records the omitted strings in fidelity receipts bound to the affected bodies.  No retained line contains a citation, so the wrapper carries no bibliography and no bibliography entry.  The retained lines contain the article's own diagram code, which the wrapper typesets with the standard tikz packages; no image file is delivered and the package declares no asset dependency.  Editorial formatting only, in addition: an AMS \texttt{amsart} preamble that restates the article's own declarations and macros used by the retained lines, namely the environments \texttt{example}, \texttt{lemma} on separate counters, together with the standard packages those lines need.  The retained environments print in this document as Example 1, Lemma 1 in order of appearance, whereas the article prints the same items with the numbering of its own full document.  The complete native source of the article as distributed in the arXiv e-print for this exact version is bundled unchanged as \texttt{sources/190231/source0.tex}.
\begin{equation}
\label{spectral-problem}
\left(\begin{array}{ll} -i \partial_x & i u\\ -i u & i \partial_x \end{array} \right) \varphi = \zeta \varphi,
\end{equation}

\begin{figure}[htb!]
	
	\caption{The pre-image of the mapping $[0,K]  \times [0,iK'] \ni \gamma \to \zeta \in [0,\infty)$ on the complex plane when $\zeta$ changes from $\zeta = 0$ to $\zeta = \infty$.}
	\label{fig-image}
\end{figure}

	\begin{equation}
	\label{Wei-Jac-3}
	H(z+iK') = i e^{\frac{\pi K'}{4 K}} e^{-\frac{i \pi z}{2K}} \Theta(z), \quad 
	\Theta(z+iK') = i e^{\frac{\pi K'}{4 K}} e^{-\frac{i \pi z}{2K}} H(z),
	\end{equation}

	\begin{equation}
	\label{exact-2}
	\left( \begin{matrix} p_1(x) \\ q_1(x) \end{matrix} \right) = 
	\left( \begin{matrix} 1 \\ -1 \end{matrix} \right) \frac{\Theta(\nu x + \alpha)}{\Theta(\nu x - \alpha)} e^{-\nu x \frac{H'(2\alpha)}{H(2\alpha)}}
	\end{equation}

	\begin{equation}
	\label{exact-1}
	\left( \begin{matrix} p_2(x) \\ q_2(x) \end{matrix} \right) = 
	\left( \begin{matrix} 1 \\ 1 \end{matrix} \right) \frac{\Theta(\nu x - \alpha)}{\Theta(\nu x + \alpha)} 
	e^{\nu x \frac{H'(2\alpha)}{H(2\alpha)}},
	\end{equation}

	\begin{equation}
	\label{elliptic-double-alpha}
	\sn(2\alpha) = \frac{\nu}{\zeta_1}, \quad \cn(2\alpha) = \frac{\zeta_3}{\zeta_1}, \quad \dn(2\alpha) = \frac{\zeta_2}{\zeta_1},
	\end{equation}

\begin{equation}
\label{Jacobi-sn}
\sn(x) = \frac{H(x)}{\sqrt{k} \Theta(x)}, \quad \cn(x) =  \frac{\sqrt{k'} H(x+K)}{\sqrt{k} \Theta(x)}, \quad \dn(x) = \frac{\sqrt{k'} \Theta(x + K)}{\Theta(x)}.
\end{equation}

\begin{equation}
\label{phinova}
\phi(x) = \frac{1}{2}\frac{\wp'(x-\frac{v}{2})+\wp'(x+\frac{v}{2})}{\wp(x-\frac{v}{2})-\wp(x + \frac{v}{2})} = \zeta\left(x+\frac{v}{2}\right)-\zeta\left(x-\frac{v}{2}\right)-\zeta(v),
\end{equation}

\begin{equation}
\label{v-alpha-correspondence}
\frac{v}{2} = -\frac{iK' + \alpha}{\sqrt{e_1-e_3}}.
\end{equation}

	\begin{equation}
	\label{Lame-disp}
	\zeta^2 = \zeta_1^2 - (\zeta_1^2 - \zeta_2^2) \sn^2(\gamma) = 
	\zeta_2^2 + (\zeta_1^2 - \zeta_2^2) \cn^2(\gamma) = 
	\zeta_3^2 + (\zeta_1^2 - \zeta_3^2) \dn^2(\gamma).
	\end{equation}

\begin{equation}
\label{Lame-sol-1}
\left( \begin{matrix} p_1(x) \\ q_1(x) \end{matrix} \right) = 
\left( \begin{matrix} 
c_1^+ \frac{ H(\nu x + \gamma + \alpha)}{\Theta(\nu x + \alpha)} + c_1^- \frac{H(\nu x + \gamma - \alpha)}{\Theta(\nu x - \alpha)} \\
c_1^+ \frac{ H(\nu x + \gamma + \alpha)}{\Theta(\nu x + \alpha)} - c_1^- \frac{H(\nu x + \gamma - \alpha)}{\Theta(\nu x - \alpha)} \end{matrix} 
\right)  e^{-\nu x Z(\gamma)}
\end{equation}

\begin{equation}
\label{Lame-sol-2}
\left( \begin{matrix} p_2(x) \\ q_2(x) \end{matrix} \right) = 
\left( \begin{matrix} c_2^+ \frac{ H(\nu x - \gamma + \alpha)}{\Theta(\nu x + \alpha)} + c_2^- \frac{H(\nu x - \gamma - \alpha)}{\Theta(\nu x - \alpha)} \\
c_2^+ \frac{ H(\nu x - \gamma + \alpha)}{\Theta(\nu x + \alpha)} - c_2^- \frac{H(\nu x - \gamma - \alpha)}{\Theta(\nu x - \alpha)}  \end{matrix} \right) 
e^{\nu x Z(\gamma)}. 
\end{equation}

\begin{example}
If $\zeta = 0$, we have $\gamma = 2 \alpha + i K'$, see Figure \ref{fig-image}.
Relations (\ref{Wei-Jac-3}) imply that 
$$
\frac{\Theta'(\gamma)}{\Theta(\gamma)} = \frac{\Theta'(2 \alpha + i K')}{\Theta(2 \alpha + i K')} = -\frac{i \pi}{2K} + \frac{H'(2\alpha)}{H(2\alpha)}.
$$
The expression (\ref{Lame-sol-1}) with $\gamma = 2 \alpha + i K'$ yields (\ref{exact-2})  if $c_1^+ = 0$ since 
$$
\frac{H(\nu x + \gamma - \alpha)}{\Theta(\nu x - \alpha)} e^{-\nu x \frac{\Theta'(\gamma)}{\Theta(\gamma)}} = 
\frac{H(\nu x + \alpha + iK')}{\Theta(\nu x - \alpha)} e^{-\nu x \frac{\Theta'(2 \alpha + i K')}{\Theta(2 \alpha + i K')}} = 
i e^{\frac{\pi K'}{4 K}} \frac{\Theta(\nu x + \alpha)}{\Theta(\nu x - \alpha)}
e^{-\nu x \frac{H'(2\alpha)}{H(2\alpha)}}.
$$
The expression (\ref{Lame-sol-2}) with $\gamma = 2 \alpha + i K'$ yields (\ref{exact-1}) if $c_2^- = 0$ since 
$$
\frac{H(\nu x - \gamma + \alpha)}{\Theta(\nu x + \alpha)} e^{\nu x \frac{\Theta'(\gamma)}{\Theta(\gamma)}} = 
\frac{H(\nu x - \alpha - iK')}{\Theta(\nu x + \alpha)} e^{\nu x \frac{\Theta'(2 \alpha + i K')}{\Theta(2 \alpha + i K')}} = 
-i e^{\frac{\pi K'}{4 K}} \frac{\Theta(\nu x - \alpha)}{\Theta(\nu x + \alpha)}
e^{\nu x \frac{H'(2\alpha)}{H(2\alpha)}}.
$$
The proportionality factors in both expressions are not relevant due to the scalar multiplication of eigenfunctions. Thus, if $\zeta = 0$, then $c_1^+ = 0$ and $c_2^- = 0$. 
\label{example-2}
\end{example}

\begin{lemma}
	\label{lem-coefficients}
If the representations (\ref{Lame-sol-1}) and (\ref{Lame-sol-2}) are solutions of the spectral problem (\ref{spectral-problem}), then their coefficients are uniquely defined up to the constant multiplication factor as 
\begin{equation}
\label{coefficient-1}
	c_1^+ = -i \zeta, \quad c_1^- = \zeta_1 \frac{\Theta(0) \Theta(2\alpha + \gamma)}{\Theta(\gamma) \Theta(2\alpha)}
\end{equation}
and
\begin{equation}
\label{coefficient-2}
c_2^+ = \zeta_1 \frac{\Theta(0) \Theta(2\alpha + \gamma)}{\Theta(\gamma) \Theta(2\alpha)}, \quad c_2^- = i \zeta.
\end{equation}
\end{lemma}

\begin{proof}
To obtain (\ref{coefficient-1}), we compute the pole contributions of the elliptic solutions of the spectral problem (\ref{spectral-problem}) at 
$$
x = \frac{v}{2} = -\frac{\alpha + i K'}{\nu}.
$$ 
Since $\wp(x) = \frac{1}{x^2} + \tilde{\wp}(x)$ as $x \to 0$ with analytic  $\tilde{\wp}(x)$ near $x = 0$, it follows from (\ref{phinova}) that 
\begin{equation}
\label{phi-asymp}
\phi(x) = -\frac{1}{x - \frac{v}{2}} + \mathcal{O}\left(x - \frac{v}{2} \right), \quad \mbox{\rm as} \;\; x \to \frac{v}{2}.
\end{equation}
By using (\ref{Jacobi-sn}) for $H(x) = \sqrt{k} \sn(x) \Theta(x)$, we
obtain the asymptotic expansion:
\begin{equation*}
H(x) = \sqrt{k} \Theta(0) x + \mathcal{O}(x^3), \quad \mbox{\rm as} \;\; x \to 0.
\end{equation*}
In view of (\ref{Wei-Jac-3}) and (\ref{v-alpha-correspondence}), this implies that 
\begin{equation}
\label{Theta-asymp}
\Theta(\nu x + \alpha) = -i \nu \sqrt{k} \Theta(0) e^{\frac{\pi K'}{4K}} e^{\frac{i \pi \nu}{2K} \left(x - \frac{v}{2} \right)} \left( x - \frac{v}{2} \right) + \mathcal{O}\left(x - \frac{v}{2} \right)^3, \quad \mbox{\rm as} \;\; x \to \frac{v}{2}.
\end{equation}
By using (\ref{Lame-sol-1}) and (\ref{Theta-asymp}), we obtain as $x \to \frac{v}{2}$:
\begin{align*}
& \frac{H(\nu x + \gamma + \alpha)}{\Theta(\nu x + \alpha)} e^{-\nu x Z(\gamma)}\\
&= \frac{i}{\nu \sqrt{k} \Theta(0)} e^{-\frac{\pi K'}{4K}} 
e^{(\alpha + iK') Z(\gamma)} e^{-\nu \left(x - \frac{v}{2} \right) \left[ Z(\gamma) + \frac{i \pi}{2K}\right]}
\frac{H\left(\nu\left(x - \frac{v}{2}\right) + \gamma - iK' \right)}{x - \frac{v}{2}} + \mathcal{O}\left(x - \frac{v}{2}\right) \\
&= \frac{1}{\nu \sqrt{k} \Theta(0)} 
e^{(\alpha + iK') Z(\gamma) + \frac{i \pi \gamma}{2K}} e^{-\nu \left(x - \frac{v}{2} \right) Z(\gamma)}
\frac{\Theta\left(\nu\left(x - \frac{v}{2}\right) + \gamma \right)}{x - \frac{v}{2}} + \mathcal{O}\left(x - \frac{v}{2}\right)
\end{align*}
and
\begin{align*}
\frac{H(\nu x + \gamma - \alpha)}{\Theta(\nu x - \alpha)} e^{-\nu x Z(\gamma)}
&= e^{(\alpha + iK') Z(\gamma)} \frac{H(\gamma - 2 \alpha - iK')}{\Theta(-2\alpha - iK')} + \mathcal{O}\left(x - \frac{v}{2}\right) \\
&= e^{(\alpha + iK') Z(\gamma) + \frac{i \pi \gamma}{2K}} \frac{\Theta(\gamma - 2 \alpha)}{H(-2\alpha)} + \mathcal{O}\left(x - \frac{v}{2}\right).
\end{align*}
In order to substitute the meromorphic parts into $p' = i \zeta p + \phi q$, 
we deduce that 
\begin{align}
p(x) &=  \frac{c_1^+ }{ \nu \sqrt{k} \Theta(0)} 
e^{(\alpha + iK') Z(\gamma) + \frac{i \pi \gamma}{2K}} 
\frac{\Theta(\gamma)}{x - \frac{v}{2}} + \mathcal{O}(1), 
\quad \mbox{\rm as} \;\; x \to \frac{v}{2},
\label{p-asymp}
\end{align}
which yields
\begin{align}
p'(x) &= -\frac{c_1^+ }{ \nu \sqrt{k} \Theta(0)} 
e^{(\alpha + iK') Z(\gamma) + \frac{i \pi \gamma}{2K}} 
\frac{\Theta(\gamma)}{\left( x - \frac{v}{2} \right)^2} 
+ \mathcal{O}(1), 
\quad \mbox{\rm as} \;\; x \to \frac{v}{2}.
\label{p-der-asymp}
\end{align}
On the other hand, we get 
\begin{align}
q(x) &= \frac{c_1^+ }{\nu \sqrt{k}\Theta(0)} 
e^{(\alpha + iK') Z(\gamma) + \frac{i \pi \gamma}{2K}} e^{-\nu \left(x - \frac{v}{2} \right) Z(\gamma)}
\frac{\Theta\left(\nu\left(x - \frac{v}{2}\right) + \gamma \right)}{x - \frac{v}{2}} \notag \\
& \quad - c_1^- e^{(\alpha + iK') Z(\gamma) + \frac{i \pi \gamma}{2K}} \frac{\Theta(\gamma - 2 \alpha)}{H(-2\alpha)} 
+ \mathcal{O}\left(x - \frac{v}{2}\right),
\quad \mbox{\rm as} \;\; x \to \frac{v}{2}.
\label{q-asymp}
\end{align}
The double pole in $p'(x)$ given by (\ref{p-der-asymp}) 
cancels out with the double pole in $\phi(x) q(x)$ given by 
the product of (\ref{phi-asymp}) and (\ref{q-asymp}). The simple pole 
from $i \zeta p(x)$ given by (\ref{p-asymp}) must cancel 
the simple pole in $\phi(x) q(x)$, due to a relation 
between $c_1^+$ and $c_1^-$. Using (\ref{phi-asymp}), (\ref{p-asymp}), (\ref{p-der-asymp}), and (\ref{q-asymp}) in $p' = i \zeta p + \phi q$, 
we obtain 
\begin{equation*}
i \zeta \frac{c_1^+ \Theta(\gamma)}{\nu \sqrt{k} \Theta(0)} -  \frac{c_1^+ }{\sqrt{k} \Theta(0)} \left[ - \Theta(\gamma) Z(\gamma) + \Theta'(\gamma) \right] + c_1^- \frac{\Theta(\gamma - 2 \alpha)}{H(-2\alpha)} = 0, 
\end{equation*}
after cancelling the constant factor $e^{(\alpha + iK') Z(\gamma) + \frac{i \pi \gamma}{2K}}$. Since $ \Theta'(\gamma) = \Theta(\gamma) Z(\gamma)$, this yields 
the linear equation 
\begin{equation}
\label{rel-c-1}
i \zeta  \frac{c_1^+ \Theta(\gamma)}{\nu \sqrt{k} \Theta(0)} 
- c_1^- \frac{\Theta(2 \alpha - \gamma)}{H(2\alpha)} = 0.
\end{equation}
The same equation follows also from $q' = -i\zeta q + \phi p$ as $x \to \frac{v}{2}$. 

Repeating the derivation as $x \to -\frac{v}{2}$, where 
$$
\left\{ \begin{array}{l} 
\phi(x) = \frac{1}{x + \frac{v}{2}} + \mathcal{O}\left(x + \frac{v}{2}\right), \\
\Theta(\nu x - \alpha) = i \nu \sqrt{k}\Theta(0)  e^{\frac{\pi K'}{4K}} e^{-\frac{i \pi \nu}{2K} \left(x + \frac{v}{2} \right)} \left( x + \frac{v}{2} \right) + \mathcal{O}\left(x + \frac{v}{2} \right)^3, \end{array} \right. \quad \mbox{\rm as} \;\; x \to -\frac{v}{2}, 
$$
we obtain another linear equation 
\begin{equation}
\label{rel-c-2}
c_1^+ \frac{\Theta(2 \alpha + \gamma)}{H(2\alpha)} + i \zeta \frac{c_1^-  \Theta(\gamma)}{\nu \sqrt{k} \Theta(0)} = 0.
\end{equation}
The determinant of the linear system (\ref{rel-c-1}) and (\ref{rel-c-2}) is zero if 
\begin{equation}
\label{det-eq-from-lin-syst}
\zeta^2 = \frac{\nu^2 k \Theta^2(0) \Theta(2\alpha-\gamma) \Theta(2\alpha + \gamma)}{\Theta^2(\gamma) H^2(2\alpha)}.
\end{equation}
Since $H(x) = \sqrt{k} \sn(x) \Theta(x)$ and 
$$
\Theta^2(0) \Theta(2\alpha-\gamma) \Theta(2\alpha+\gamma) = \Theta^2(2\alpha) \Theta^2(\gamma) - H^2(2\alpha) H^2(\gamma), 
$$
we use the relation $\nu = \zeta_1 \sn(2\alpha)$ from (\ref{elliptic-double-alpha}) and obtain from (\ref{det-eq-from-lin-syst}):
\begin{align*}
\zeta^2 &= \nu^2 k \left[ \frac{\Theta^2(2 \alpha)}{H^2(2 \alpha)} - \frac{H^2(\gamma)}{\Theta^2(\gamma)} \right] \\
&= \nu^2 \left[ \frac{1}{\sn^2(2\alpha)} - k^2 \sn^2(\gamma) \right] \\
&= \zeta_1^2 -  (\zeta_1^2 - \zeta_2^2) \sn^2(\gamma),
\end{align*}
which recovers the dispersion relation (\ref{Lame-disp}). Therefore, one 
of the two equations in the linear system (\ref{rel-c-1}) and (\ref{rel-c-2}) is redundant.

Recall from Example \ref{example-2} that $\gamma = 2 \alpha + i K'$ for $\zeta = 0$, which yields $c_1^+ = 0$ since $\Theta(2\alpha - \gamma) = 0$ and $\Theta(2\alpha + \gamma) \neq 0$. Therefore, we use (\ref{rel-c-2}) of the linear system (\ref{rel-c-1}) and (\ref{rel-c-2}), set $c_1^+ = -i\zeta$ and obtain 
$$
c_1^- = \frac{\nu \sqrt{k} \Theta(0) \Theta(2\alpha + \gamma)}{\Theta(\gamma) H(2\alpha)} = \zeta_1 \frac{\Theta(0) \Theta(2\alpha + \gamma)}{\Theta(\gamma) \Theta(2\alpha)},
$$
where we have used again that $\nu = \zeta_1 \sn(2\alpha)$. This yields (\ref{coefficient-1}).

To obtain (\ref{coefficient-2}) with a similar method, we expand the second solution (\ref{Lame-sol-2}) as $x \to \pm \frac{v}{2}$.  Here, we use the fact that (\ref{Lame-sol-2}) 
is obtained from (\ref{Lame-sol-1}) by replacing $\gamma$ to $-\gamma$. Since the poles $\pm \frac{v}{2}$ are independent of $\gamma$, we obtain the linear system of equations for $(c_2^+,c_2^-)$ by replacing $\gamma$ to $-\gamma$ in (\ref{rel-c-1}) and (\ref{rel-c-2}):
\begin{equation}
\label{rel-c-3}
i \zeta  \frac{c_2^+ \Theta(\gamma)}{\nu \sqrt{k} \Theta(0)} 
- c_2^- \frac{\Theta(2 \alpha + \gamma)}{H(2\alpha)} = 0
\end{equation}
and
\begin{equation}
\label{rel-c-4}
c_2^+ \frac{\Theta(2 \alpha - \gamma)}{H(2\alpha)} + i \zeta \frac{c_2^-  \Theta(\gamma)}{\nu \sqrt{k} \Theta(0)} = 0.
\end{equation}
Again, it follows from Example \ref{example-2} for $\zeta = 0$ that $c_2^- = 0$ since $\Theta(2\alpha - \gamma) = 0$ and $\Theta(2\alpha + \gamma) \neq 0$. 
Therefore, we use (\ref{rel-c-3}) of the linear system 
(\ref{rel-c-3}) and (\ref{rel-c-4}), set $c_2^- = i\zeta$ and obtain 
$$
c_2^+ = \frac{\nu \sqrt{k} \Theta(0) \Theta(2\alpha + \gamma)}{\Theta(\gamma) H(2\alpha)} = \zeta_1 \frac{\Theta(0) \Theta(2\alpha + \gamma)}{\Theta(\gamma) \Theta(2\alpha)}.
$$
This yields (\ref{coefficient-2}).
\end{proof}

\end{document}
